% Apéndice: rotaciones infinitesimales y orden de composición
\section{Rotaciones infinitesimales y orden de composición}\label{cap:15}\label{ap:infinitesimales}

\subsection{De dónde surge \(R(\varepsilon)\approx I+\varepsilon A\)}

La aproximación surge del desarrollo de Taylor de la matriz de rotación alrededor del ángulo
cero. Para una rotación de ángulo \(\varepsilon\) alrededor del eje \(x\):
\[
R_x(\varepsilon)=
\begin{pmatrix}
1&0&0\\
0&\cos\varepsilon&-\sin\varepsilon\\
0&\sin\varepsilon& \cos\varepsilon
\end{pmatrix}.
\]
Las series de Taylor son
\[
\cos\varepsilon
=1-\frac{\varepsilon^2}{2!}+\frac{\varepsilon^4}{4!}-\cdots,
\qquad
\sin\varepsilon
=\varepsilon-\frac{\varepsilon^3}{3!}+\frac{\varepsilon^5}{5!}-\cdots.
\]
Sustituyendo:
\[
R_x(\varepsilon)=
\begin{pmatrix}
1&0&0\\
0&1-\frac{\varepsilon^2}{2}+\cdots&
-\varepsilon+\frac{\varepsilon^3}{6}-\cdots\\
0&\varepsilon-\frac{\varepsilon^3}{6}+\cdots&
1-\frac{\varepsilon^2}{2}+\cdots
\end{pmatrix}.
\]
Si se conservan únicamente los términos constantes y lineales en \(\varepsilon\):
\[
R_x(\varepsilon)
=
\underbrace{
\begin{pmatrix}
1&0&0\\0&1&0\\0&0&1
\end{pmatrix}}_{I}
+
\varepsilon
\underbrace{
\begin{pmatrix}
0&0&0\\0&0&-1\\0&1&0
\end{pmatrix}}_{A_x}
+O(\varepsilon^2).
\]
Por lo tanto,
\[
\boxed{R_x(\varepsilon)=I+\varepsilon A_x+O(\varepsilon^2)}.
\]
El símbolo \(O(\varepsilon^2)\) reúne todos los términos cuyo tamaño disminuye como
\(\varepsilon^2\) o más rápidamente cuando \(\varepsilon\to0\).

\begin{resultado}[Forma general]
Para una familia suave de rotaciones basada en la identidad, es decir, con \(R(0)=I\), Taylor
alrededor de \(\theta=0\) da
\[
R(\varepsilon)=R(0)+\varepsilon R'(0)+O(\varepsilon^2).
\]
Como \(R(0)=I\), definiendo \(A=R'(0)\) se obtiene
\[
\boxed{R(\varepsilon)=I+\varepsilon A+O(\varepsilon^2)}.
\]
La matriz \(A\) se denomina \textbf{generador infinitesimal} de la rotación.
\end{resultado}

\subsection{Generadores alrededor de los tres ejes}

Con la convención de giros positivos de la mano derecha:
\[
A_x=
\begin{pmatrix}
0&0&0\\
0&0&-1\\
0&1&0
\end{pmatrix},
\quad
A_y=
\begin{pmatrix}
0&0&1\\
0&0&0\\
-1&0&0
\end{pmatrix},
\quad
A_z=
\begin{pmatrix}
0&-1&0\\
1&0&0\\
0&0&0
\end{pmatrix}.
\]
Estas matrices son antisimétricas, \(A_i^T=-A_i\), y describen la velocidad instantánea
con la que se mueve un vector al comenzar la rotación.

\subsection{Interpretación geométrica}

Sea \(e_x=(1,0,0)^T\). Para cualquier vector \(v\),
\[
  A_xv=e_x\times v.
\]
Entonces:
\[
  R_x(\varepsilon)v\approx v+\varepsilon(e_x\times v).
\]
El primer término es el vector original. El segundo es un pequeño desplazamiento perpendicular
al eje \(e_x\) y al vector \(v\), es decir, tangente a la circunferencia recorrida durante
la rotación.

\begin{center}
\begin{tikzpicture}[scale=1.0]
  \coordinate (o) at (1,1);
  \coordinate (v) at (6.2,2.0);
  \coordinate (vp) at (5.85,2.95);
  \draw[axisline,-{Latex[length=2.5mm]}] (o)--(8,1)
    node[right] {eje \(x\)};
  \draw[mainline,-{Latex[length=2.8mm]}] (o)--(v)
    node[midway,below] {$v$};
  \draw[rust,line width=1.05pt,-{Latex[length=2.8mm]}] (v)--(vp)
    node[right] {$\varepsilon A_xv$};
  \draw[teal,dashed] (o)--(vp)
    node[midway,above left] {$R_x(\varepsilon)v$};
  \fill[ink] (o) circle (2.4pt);
  \fill[teal] (v) circle (2.4pt);
  \fill[rust] (vp) circle (2.4pt);
  \node[font=\small,text=muted,align=center] at (5.4,3.8)
    {el desplazamiento infinitesimal\\es tangente a la trayectoria};
\end{tikzpicture}
\end{center}

\begin{idea}[Qué significa geométricamente la aproximación]
Una rotación finita sigue un arco de circunferencia. En un intervalo angular muy pequeño,
el arco y su recta tangente son indistinguibles a primer orden. La expresión
\(v+\varepsilon A_xv\) reemplaza localmente el arco por ese desplazamiento tangencial.
\end{idea}

\subsection{Por qué el orden no aparece a primer orden}

Consideremos dos rotaciones pequeñas:
\[
R_x(\varepsilon)=I+\varepsilon A_x+O(\varepsilon^2),
\qquad
R_y(\delta)=I+\delta A_y+O(\delta^2).
\]
Primero \(R_y\) y después \(R_x\):
\[
\begin{aligned}
R_x(\varepsilon)R_y(\delta)
&\approx(I+\varepsilon A_x)(I+\delta A_y)\\
&=I+\varepsilon A_x+\delta A_y+\varepsilon\delta A_xA_y.
\end{aligned}
\]
En el orden inverso:
\[
\begin{aligned}
R_y(\delta)R_x(\varepsilon)
&\approx(I+\delta A_y)(I+\varepsilon A_x)\\
&=I+\varepsilon A_x+\delta A_y+\varepsilon\delta A_yA_x.
\end{aligned}
\]
Los términos de primer orden son exactamente los mismos. La diferencia es
\[
\boxed{
R_x(\varepsilon)R_y(\delta)
-R_y(\delta)R_x(\varepsilon)
\approx\varepsilon\delta[A_x,A_y]
},
\]
donde
\[
  [A_x,A_y]=A_xA_y-A_yA_x
\]
es el conmutador.

\begin{resultado}[Conclusión precisa]
Las rotaciones infinitesimales no conmutan exactamente. Parecen conmutar cuando se trabaja
solo a primer orden porque la diferencia causada por el orden es proporcional al producto
\(\varepsilon\delta\), que es de segundo orden.
\end{resultado}

\subsection{Interpretación geométrica del cambio de orden}

La primera rotación modifica ligeramente la orientación del objeto sobre el cual actúa la
segunda. Esa modificación es del orden de \(\varepsilon\). Como la segunda rotación tiene
tamaño \(\delta\), el efecto adicional producido por aplicar una sobre la orientación creada
por la otra tiene tamaño \(\varepsilon\delta\).

Si \(\varepsilon=\delta=0.01\) radianes:
\[
\varepsilon=10^{-2},
\qquad
\varepsilon\delta=10^{-4}.
\]
El movimiento principal tiene tamaño aproximado \(10^{-2}\), mientras que la diferencia entre
los dos órdenes aparece aproximadamente en \(10^{-4}\). Al conservar solamente la primera
cifra de pequeñez, esa diferencia desaparece.

\subsection{Relación con la exponencial matricial}

Las rotaciones pueden escribirse exactamente como
\[
  R_x(\varepsilon)=e^{\varepsilon A_x},
  \qquad
  R_y(\delta)=e^{\delta A_y}.
\]
La fórmula de Baker--Campbell--Hausdorff comienza como
\[
e^{\varepsilon A_x}e^{\delta A_y}
=
\exp\left(
\varepsilon A_x+\delta A_y
+\frac12\varepsilon\delta[A_x,A_y]+\cdots
\right).
\]
El conmutador aparece explícitamente como la primera corrección que registra el orden.
Para los generadores de rotaciones:
\[
  [A_x,A_y]=A_z,
\]
de modo que la diferencia entre girar infinitesimalmente alrededor de \(x\) y de \(y\) en
órdenes opuestos se manifiesta como una pequeña rotación residual alrededor de \(z\).

\begin{atencion}[Idea que conviene retener]
\[
\boxed{
\text{A primer orden las rotaciones pequeñas se suman;}
\quad
\text{a segundo orden aparece el orden.}
}
\]
No se afirma que las rotaciones tridimensionales conmuten, sino que su no conmutatividad es
demasiado pequeña para aparecer en la aproximación lineal.
\end{atencion}

